\sscp{Nasals}{13-01-01}
We look first at the behaviour of nasals.
\trf{13-01} shows the nasals in initial,
medial, and final positions for each subgroup.
We note that confident reconstructions of PMP *ñ are not common,
such that finding sufficient reflexes in any given language to
establish reliable sound correspondences is challenging.
PMP *ñ is not found in word-final position.

\begin{table}
\caption[Summary of nasals in each subgroup]{Summary of nasals in each subgroup\su{†}}
\label{tab:13-01}
\begin{adjustbox}{width=\textwidth}
	\begin{threeparttable}
	\st{2pt}\begin{tabular}{llll|lll|ll|lll}\lsptoprule
PMP					&	*{\#}m-			&	*-m-		&	*m{\#}				&	*{\#}n-	&	*-n-		&	*n{\#}				&	*{\#}ñ			&	*-ñ-						&	*{\#}ŋ			&	*-ŋ-			&	*ŋ{\#}\\	\midrule
\tsc{BL}		&	{\hsh}m			&	{\hd}m	&	{\hr}m				&	{\hsh}n	&	{\hd}n	&	{\hr}n, ŋ			&	{\hsh}ñ			&	{\hd}ñ					&	{\hsh}ŋ			&	{\hd}ŋ		&	{\hr}ŋ\\
\tsc{TB}		&	{\hsh}m			&	{\hd}m	&	{\hr}m, {\0}	&	{\hsh}n	&	{\hd}n	&	{\hr}n, {\0}	&	{\hsh}n			&	{\hd}n\tsc{+fr}	&	{\hsh}ŋ, n	&	{\hd}ŋ, n	&	{\hr}n, {\0}\\
\tsc{CT}		&	{\hsh}m			&	{\hd}m	&	{\hr}m, {\0}	&	{\hsh}n	&	{\hd}n	&	{\hr}n, {\0}	&	{\hsh}n			&	{\hd}n					&	{\hsh}g			&	{\hd}g		&	{\hr}g\\
\tsc{Aru}		&	{\hsh}m			&	{\hd}m	&	{\hr}m				&	{\hsh}n	&	{\hd}n	&	{\hr}n				&	{\hsh}ñ?		&	{\hd}n\tsc{+fr}	&	{\hsh}ŋ			&	{\hd}ŋ		&	{\hr}ŋ?\\
STB					&	{\hsh}m			&	{\hd}m	&	{\hr}m				&	{\hsh}n	&	{\hd}n	&	{\hr}n, {\0}	&	{\hsh}n			&	{\hd}n					&	{\hsh}ŋ			&	{\hd}ŋ		&	{\hr}n\\
%Banda\su{‡}	&	{\hsh}m, mb	&	{\hd}m	&	{\hr}n				&	{\hsh}n	&	{\hd}n	&	{\hr}n				&	{\hsh}n, ñ	&	{\hd}n					&	{\hsh}n			&	{\hd}n		&	{\hr}n\\
Banda				&	{\hsh}m, mb	&	{\hd}m	&	{\hr}n				&	{\hsh}n	&	{\hd}n	&	{\hr}n				&	{\hsh}n, ñ	&	{\hd}n					&	{\hsh}n			&	{\hd}n		&	{\hr}n\\
\tsc{AS}		&	{\hsh}m			&	{\hd}m	&	{\hr}m, {\0}	&	{\hsh}n	&	{\hd}n	&	{\hr}n, {\0}	&	{\hsh}n			&	{\hd}n\tsc{+fr}	&	{\hsh}n			&	{\hd}n		&	{\hr}n, {\0}\\
\tsc{SB}		&	{\hsh}m			&	{\hd}m	&	{\hr}n, {\0}	&	{\hsh}n	&	{\hd}n	&	{\hr}n, {\0}	&	{\hsh}n			&	{\hd}n\tsc{+fr}	&	{\hsh}ŋ			&	{\hd}ŋ		&	{\hr}n, {\0}\\
		\lspbottomrule
	\end{tabular}
		\begin{tablenotes}\tnsize
			\item[†] {\tsc{+fr} = fronting of a preceding or following vowel as traces
						of *ñ in some words.
						Technically in this and following tables, each value should be
						preceded by * indicating it is the proto-phoneme for that subgroup,
						but we avoid doing so to reduce visual clutter
						and for the sake of space.
						A question mark (?) indicates there is not enough data available
						to have confidence in the correspondence.
						See \prf{sec:Subgroups} for subgroup abbreviations.}
			%\item[‡]	{Banda has initial and medial [ɲ] in some forms that are nasal
								%allophones assimilated to point of articulation of the following C.}
		\end{tablenotes}
	\end{threeparttable}
\end{adjustbox}
\end{table}

The data show that all distinctions must be reconstructed for
the region as a whole.
No splits or mergers are shared by all subgroups.
The one exception is the apparent merger of *n/*ñ > \it{n} in several subgroups.
However, we note:

\begin{enumerate}
	\item the forms reconstructed with PMP *ñ frequently
				change reflecting the uncertain status of the proto-phoneme,
				such that we sometimes encounter *ñiuR {\tl} *niuR ‘coconut’,
				or *buñi {\tl} *buni ‘hide,
				conceal’ leaving a moving target that is often hard to pin
				down with confidence;
	\item the data in any individual language
				or group of languages for *ñ are quite sparse;
	\item the merger
				of *n/*ñ > \it{n} is very common,
				which lessens its value as an innovation for subgrouping purposes
				and increases the likelihood of parallel development.
\end{enumerate}

For example, the same merger also occurs in contiguous subgroups
just west of the Blust line across the entire \tsc{Bungku-Tolaki} subgroup
(\citealp[55ff]{Mead1998}; \citealp{Mead1999}),
so the merger of *n/*ñ > \it{n} is not exclusive to parts
of the target region.
Consequently we have little basis or confidence for claiming
this apparent merger is an exclusively shared
innovation with relevance for subgrouping.
At best, it is weak.

Furthermore, for example, in \tsc{Timor-Babar},
PMP *ñ often triggers raising of a following vowel.
Examples include the following (for all these examples final
*a/*ə would be expected to become \it{a}: PMP *utaña ‘ask’ >
Tii, Lole, Termanu (Rote) \it{na-tan\tbr{e}}, Kisar \it{-aʔukan\tbr{i}}
‘ask’ (*t > \it{k} is regular),
but Luang \it{hoi-tian\tbr{a}}; PMP *añam/*añəm ‘plait,
braid’ > Helong \it{nan\tbr{e}ŋ}
and Molo \it{an\tbr{i}} ‘braid’; P(W)MP *bañən ‘sneeze’ > Termanu
\it{kisu-fan\tbr{i}}, Dela \it{ʔesu-fan\tbr{i}}, East Tetun \it{fan\tbr{i}} ‘sneeze’.
Such examples indicate that it may be necessary to posit that
*ñ was retained unchanged at the level of Proto-Timor-Babar
and thus caution against using the merger of *ñ/*n as a defining
sound change for a larger subgroup containing \tsc{Timor-Babar} languages.
Similar traces of PMP *ñ in surrounding vowels are found in other subgroups.

\citet{Schapper2015} noted that many languages in our target
region have no engma /ŋ/.
She was not claiming that the absence of engma had subgrouping
value, but rather observing a fairly widespread feature in the region
in both Austronesian (circle)
and Papuan (square) languages.
\frf{13-01} shows the distribution of languages with a known absence of engma.
White circles indicate languages that were investigated,
but retain engma in their native vocabulary.
Data for the map is drawn from \citet{DonohueDawson2013},
supplemented by around 50 languages from our own sources (see appendix \ref{ch:LanLis}).

Many languages retain PMP *ŋ = \it{ŋ}.
Some have lost the phoneme by merging PMP *ŋ/*n > \it{n},
yet do have an engma in loans or through other processes.
For example, \tsc{Ambon-Seram} language Asilulu merges PMP *ŋ/*n
> \it{n}, but nevertheless Asilulu has many words with \it{ŋ}.
So PMP *\tbr{ŋ}ajan > \it{\tbr{n}ala-} ‘name’; PMP *na\tbr{ŋ}uy >
\it{na\tbr{n}u} ‘swim’; PMP *də\tbr{ŋ}əR > \it{pa-ne\tbr{n}e}
‘hear, listen’; but \it{baŋka} ‘swollen’ < Ambon Malay \it{baŋka} (cf.
Indonesian \it{bəŋkak} ‘swollen’), \it{maŋako} ‘agree’ < Indonesian
\it{məŋaku} ‘confess, acknowledge’, \it{ŋaʤi} ‘recite the Quran’ < Indonesian \it{məŋaʤi
{\tl} məŋkaʤi} ‘recite the Quran’,
\it{oyaŋ} ‘great grandparent’ < Malay \it{moyaŋ} ‘great grandparent, ancestor, forefather’.
Many \tsc{Ambon-Seram} languages have the same pattern as Asilulu:
no retentions of PMP *ŋ¸ but \it{ŋ} is found in many loanwords.
In Makatian, *ŋ > \it{n} as in \it{la\tbr{n}it-e} ‘sky’ (< *la\tbr{ŋ}it),
\it{mri\tbr{n}in} ‘cold’ (< *\it{ma-di\tbr{ŋ}diŋ}),
but /n/ → [ŋ] /{\gap}k,
as in \it{ŋkaye} ‘boat’, \it{ŋkat} ‘bite’.
In the figure below,
Makatian would be treated as not having \it{ŋ}.
Given the severe under-documentation of languages in the region,
the presence or absence of engma is not always clear.

Notice that large pockets of languages
and whole subgroups retain engma,
such as \tsc{Sula-Buru}, \tsc{Aru},
nodes of \tsc{Seram-Tanimbar-Bomberai}, and \tsc{Bima-Lembata} languages.
All distinctions in the nasals must be reconstructed to account
for the data across Wallacean subgroups.

\begin{figure}[h]\centering
	\includegraphics[width=\columnwidth]{engma.pdf}
	\fcp{Distribution of absence of engma /ŋ/}{13-01}
\end{figure}
